\section{Discussion}\label{sec:discussion}

\subsection{KG Infrastructure as Operational Backbone, Not LLM Interface}

Our results suggest a reframing of KG infrastructure value in the
LLM era. The conventional motivation for deploying a hyper-relational
KG in front of an LLM is the expressivity argument: the KG can
faithfully represent N-ary facts and structural constraints that
binary triples or natural language cannot. Our data shows that this
expressivity argument does not transfer to the LLM-facing
interface — given the same facts in-context, all surface forms
including natural language reach $\geq 83\%$ accuracy, with NL the
Pareto winner on accuracy and token cost simultaneously.

The value of KG infrastructure therefore lies elsewhere. We identify
four operational dimensions in which KGs remain critical:

\begin{enumerate}
\item \textbf{Production scalability.} Schema enforcement, integrity
      constraints, concurrency control, and ACID transactions are
      properties of the database backbone, not of any single
      LLM-facing prompt. A KG store provides these for free; an
      LLM-prompt pipeline that holds facts only in JSON files has to
      re-implement them.
\item \textbf{Multi-system source of truth.} A KG can serve a chat
      LLM, a BI dashboard, an audit pipeline, and an internal API
      from a single fact store. The LLM interface is one of many
      consumers; the KG's role is integration substrate.
\item \textbf{Boundary-and-rule encoding.} Schema constraints
      (cardinality, role types, mandatory attributes) and explicit
      negation queries (\texttt{not \{...\};}) let a KG agent reason
      about \emph{absence} in a way an in-context dump cannot. The
      mechanism findings of Section~\ref{sec:mechanism} suggest that
      the largest residual KG advantage — closed-world boundary
      inference — is best exploited via the agent's query interface,
      not the LLM's reading of a raw dump.
\item \textbf{Latent expressivity for stronger query agents.} Today
      the query-generation LLM is the bottleneck. As specialized
      coder LLMs improve at TypeQL, Cypher, and similar domain
      languages, the production KG agent's accuracy is bounded above
      by KG expressivity, not by surface form. The current 30-point
      gap may be closeable from the agent side, not the
      surface-form side.
\end{enumerate}

\subsection{Practical Recommendations for KG-LLM Pipelines}

Our findings yield four concrete recommendations.

\paragraph{Recommendation 1: Add a KG-to-NL conversion layer.} For
in-context QA, structure-preserving NL summary outperforms TypeQL,
Cypher, and JSON-LD on both accuracy and token cost. Generation can
be templated (one sentence per fact) and is fully deterministic
from the source data. Production should ship this layer between the
KG store and the chat LLM, treating the structured surface forms as
debug or audit views rather than the LLM-facing form.

\paragraph{Recommendation 2: Augment NL with explicit boundary
annotations.} The mechanism analysis showed that even NL loses some
out-of-scope accuracy because the dump enumerates only positive
facts. Adding a boundary clause per relation — \emph{"X is defined
only for grades $a$ and $b$; other grades are not in this table"} —
is a deterministic transformation from the schema and should lift
out-of-scope accuracy further. We treat measurement of this lift
as future work.

\paragraph{Recommendation 3: Specialize the query-generation LLM.}
Setting B's query-model fix to Qwen3-Coder partially explains why
Llama-4-Scout's KG agent (56\%) is worse than its in-context tier
(63\%): the chat-side syntax weakness compounds with the absent
syntax-aware specialization on the query side. As the available
specialized coder models grow more competent at KG query languages,
we expect the production-agent / in-context gap to narrow from the
agent side. Practitioners should track coder-LLM benchmarks for
TypeQL and Cypher as a leading indicator.

\paragraph{Recommendation 4: Treat KG as backbone, not interface.}
The four dimensions in Section~7.1 motivate keeping the KG store as
the source of truth and the integration substrate. The LLM
interface above it should be a thin NL-summary layer with optional
agent fall-through for queries the summary cannot answer (e.g.,
counts, filters, multi-hop traversals beyond the summary template).
This inverts the common architecture in which the KG agent is
\emph{primary} and a context document is the fall-back.

\subsection{Scope: Fact-Bounded KGs in the 32K-Token Regime}

A natural objection is that our findings apply only when the entire
KG fits in the LLM context, and thus cannot generalize to KGs serving
Web-scale or multi-million-fact use cases. We accept this scope but
argue the relevant population is larger than it first appears.
Regulations, internal policies, benefits schedules, clinical
guidelines, organizational handbooks, compliance frameworks, and many
enterprise knowledge bases serialize to a few megabytes of structured
facts — well inside the 32K--200K token windows of current LLMs. For
\emph{this} substantial class of domains, our results are directly
applicable: production teams currently invest in KG-agent
infrastructure where the cheaper, more accurate path is an in-context
NL summary. The argument for KG infrastructure shifts from "in-context
limits" (where it does not hold for fact-bounded KGs) to the four
operational dimensions in \S7.1. For Web-scale or open-domain KGs, the
retrieval problem is different and our findings do not extend; we
treat this as an explicit boundary.

\subsection{Relation to Counterintuitive H1}

The pre-experiment hypothesis — "even a maximally preprocessed
context-RAG cannot match a hyper-relational KG agent on multi-key
queries" — is rejected by our data ($p<0.05$ in 5/8 KG-vs-context
comparisons, all in the opposite direction). The data
\emph{does} support a different counterintuitive claim: that the
production KG agent's query-generation stage, not the KG's
expressivity, is the dominant accuracy determinant in this
deployment regime. This new claim shares the original's
counterintuitive framing — practitioners typically assume that more
structure leads to better LLM downstream performance — while being
directly evidenced by the surface-form ablation.
